% Generated draft: requires mathematical and source review.
\documentclass[11pt]{article}
\usepackage[margin=25mm]{geometry}
\usepackage{fontspec}
\setmainfont{Latin Modern Roman}
\usepackage{amsmath,amssymb,mathtools,unicode-math}
\setmathfont{Latin Modern Math}
\usepackage{enumitem,needspace,xurl,hyperref}
\setlength{\parindent}{0pt}
\setlength{\parskip}{5pt}
\setlength{\emergencystretch}{3em}
\newcommand{\N}{\mathbb{N}}
\newcommand{\Z}{\mathbb{Z}}
\newcommand{\Q}{\mathbb{Q}}
\newcommand{\R}{\mathbb{R}}
\newcommand{\C}{\mathbb{C}}
\newcommand{\F}{\mathbb{F}}
\newcommand{\A}{\mathbb{A}}
\newcommand{\PP}{\mathbb{P}}
\newcommand{\E}{\mathbb{E}}
\newcommand{\eps}{\varepsilon}
\newcommand{\dd}{\,\mathrm{d}}
\begin{document}
\section*{UnsolvedMath ID: 10400070}
\textbf{Source problem code:} catalogue entry
\subsection{Definitions and mathematical statement}
For a knot $K$ in the source's category, let $\Delta_K(t)$ be its Alexander polynomial and put
\[
R_K=\Q[t,t^{-1},\Delta_K(t)^{-1}].
\]
The rational invariant is denoted
\[
Z^{\mathrm{rat}}(K)\in\mathcal A_{R_K}(\varnothing;\Q).
\]
Its target is the source's rationally beaded diagram space with no external skeleton/legs, with the prescribed relations and completion. Edge labels lie in the stated localization, rather than arbitrary rational functions or only Laurent polynomials.

For every allowed loop degree $\ell$, the problem asks for a topological construction of the corresponding component $\pi_\ell Z^{\mathrm{rat}}(K)$. On a connected diagram, loop degree ordinarily means its first Betti number $b_1=|E|-|V|+1$. The source does not specify how to extract the loop part from a group-like series, in particular whether it means the connected/primitive part or a total-degree convention on disconnected diagrams. No convention is silently selected.

A construction must identify the specified invariant component from topological data of $K$ and establish agreement, not merely rename that component. The record does not prescribe the allowable topological operations, target normalization, orientation/bead/holonomy relations, or exact knot ambient/framing category. It also does not give a finiteness/computability requirement for ``construction''. Thus the ring and all-loop-degree request are precise as printed, while the full rational-invariant and component definitions must be recovered from the preceding source.
\subsection{Short English statement}
Construct each loop-degree component of the rational knot invariant directly from topology, with agreement in its Alexander-localized beaded diagram space. The loop extraction and full target conventions are not supplied.
\subsection{Sources}
\begin{itemize}
\item Local UnsolvedMath registry, permanent ID 10400070, exact source code catalogue entry, Problem 3.20; no proposer/year supplied. Registry status open, classification OPEN-TRIAGE.
\item Ohtsuki, \emph{Problems on invariants of knots and 3-manifolds} (2002), Problem 3.20, PDF page 70, \url{https://msp.org/gtm/2002/04/gtm-2002-04-024s.pdf}. Supplied pdf-font-structure extraction accessed July 29, 2026; provenance distinguishes 231 bold numbered items. No PDF was accessed here.
\item The supplied summary calls this the topological-construction question for rational $Z$ and says it is not obviously resolved, without a precise later reference or missing definition.
\item Provenance flags current primary-literature status and redistribution terms NEEDS\_REVIEW.
\end{itemize}
\end{document}
